\notechapter{N-20250926}{Backpropagation in One MLP Layer: From a Three-Dimensional Tensor to an Outer Product}

\section{The setting}

Consider one middle layer of a multilayer perceptron:
\[
  z=Wx+b,\qquad y=\sigma(z).
\]
Here
\[
  x\in\mathbb R^n,\qquad
  W\in\mathbb R^{m\times n},\qquad
  b\in\mathbb R^m,\qquad
  z,y\in\mathbb R^m.
\]
The activation function $\sigma$ acts elementwise, for example sigmoid or ReLU.  The goal of backpropagation is to compute
\[
  \frac{\partial L}{\partial W}
\]
for a loss function $L$.

The softmax--cross-entropy setting provides a standard example.  If the final layer produces logits $z_i$, then
\[
  \hat y_i=\frac{e^{z_i}}{\sum_j e^{z_j}},\qquad
  L=-\sum_i t_i\ln \hat y_i,
\]
where $t_i$ is a one-hot label.  Backpropagation asks how this scalar loss changes when each weight $W_{j,k}$ changes.

\section{The theoretical tensor}

Strictly speaking, the derivative of the vector $y\in\mathbb R^m$ with respect to the matrix $W\in\mathbb R^{m\times n}$ has three indices:
\[
  \frac{\partial y_i}{\partial W_{j,k}}.
\]
The index $i$ chooses the output component, $j$ chooses the row of $W$, and $k$ chooses the column of $W$.  So the object has shape
\[
  m\times m\times n.
\]

This is the source of the confusion recorded in the note.  The derivative is theoretically a third-order tensor, but in neural-network code one usually sees only a matrix gradient.  The matrix appears after chain-rule contraction with the upstream gradient.

\section{Computing the sparse derivative}

Write the linear preactivation component as
\[
  z_r=\sum_{t=1}^n W_{r,t}x_t+b_r.
\]
Then
\[
  \frac{\partial z_r}{\partial W_{j,k}}
  =
  \begin{cases}
  x_k,&r=j,\\
  0,&r\ne j,
  \end{cases}
  =\delta_{rj}x_k.
\]
This is sparse because $z_r$ depends only on row $r$ of $W$.

For $y_i=\sigma(z_i)$ with elementwise activation,
\[
  \frac{\partial y_i}{\partial z_r}=0\quad(i\ne r),\qquad
  \frac{\partial y_i}{\partial z_i}=\sigma'(z_i).
\]
Thus
\[
  \frac{\partial y_i}{\partial W_{j,k}}
  =\sum_{r=1}^m
  \frac{\partial y_i}{\partial z_r}
  \frac{\partial z_r}{\partial W_{j,k}}
  =\frac{\partial y_i}{\partial z_j}x_k.
\]
For elementwise activation this is nonzero only when $i=j$.

\section{Contraction by the chain rule}

The loss $L$ is scalar.  Therefore
\[
  \frac{\partial L}{\partial W_{j,k}}
  =\sum_{i=1}^m
  \frac{\partial L}{\partial y_i}
  \frac{\partial y_i}{\partial W_{j,k}}.
\]
Substituting the previous formula gives
\[
  \frac{\partial L}{\partial W_{j,k}}
  =\sum_{i=1}^m
  \frac{\partial L}{\partial y_i}
  \frac{\partial y_i}{\partial z_j}x_k.
\]
The sum over $i$ is exactly the $j$-th component of the upstream gradient with respect to $z$:
\[
  \delta_j=\frac{\partial L}{\partial z_j}.
\]
Hence
\[
  \frac{\partial L}{\partial W_{j,k}}=\delta_jx_k.
\]
Stacking all $j,k$ gives the outer product
\[
  \boxed{\displaystyle
  \frac{\partial L}{\partial W}=\delta x^T
  }
\]
where
\[
  \delta=\frac{\partial L}{\partial z}.
\]

\section{Why the tensor disappears}

The original derivative $\partial y/\partial W$ is a tensor with shape $(m,m,n)$.  But backpropagation does not need to store it.  The chain rule contracts it with the vector $\partial L/\partial y$, leaving a matrix of shape $(m,n)$.

Equivalently:
\[
  \text{three-dimensional tensor}
  \quad\xrightarrow{\text{chain-rule summation}}
  \quad
  \text{matrix gradient}.
\]
This is not a loss of information for gradient descent, because gradient descent only needs the derivative of the scalar loss with respect to each parameter.

\section{Contravariant reading of the same calculation}

The same computation can be described as a pullback of covectors.  A forward layer is a map
\[
  f:\text{parameters and activations}\longrightarrow \text{outputs}.
\]
Its differential sends small forward perturbations to small output perturbations:
\[
  df:T_xX\to T_{f(x)}Y.
\]
But the loss supplies a covector at the output,
\[
  dL\in T_{f(x)}^*Y,
\]
and backpropagation sends it backward by precomposition:
\[
  f^*(dL)=dL\circ df\in T_x^*X.
\]
This is why reverse-mode differentiation is naturally contravariant.

For a batched linear layer in PyTorch convention,
\[
  X\in\mathbb R^{b\times n},\qquad
  W\in\mathbb R^{m\times n},\qquad
  Y=XW^T\in\mathbb R^{b\times m},
\]
the component formula is
\[
  Y_{\alpha i}=\sum_{k=1}^n X_{\alpha k}W_{ik}.
\]
Hence
\[
  \frac{\partial Y_{\alpha i}}{\partial W_{jk}}
  =
  \delta_{ij}X_{\alpha k},
\]
a Jacobian tensor of shape \((b,m,m,n)\).  Contracting with the upstream gradient \(G=\partial L/\partial Y\), of shape \((b,m)\), gives
\[
\begin{aligned}
  \frac{\partial L}{\partial W_{jk}}
  &=
  \sum_{\alpha,i}G_{\alpha i}
  \frac{\partial Y_{\alpha i}}{\partial W_{jk}}\\
  &=
  \sum_{\alpha}G_{\alpha j}X_{\alpha k},
\end{aligned}
\]
so
\[
  \frac{\partial L}{\partial W}=G^TX.
\]
If the weights are stored instead as \(W_{\mathrm{col}}\in\mathbb R^{n\times m}\) and \(Y=XW_{\mathrm{col}}\), the same formula appears as
\[
  \frac{\partial L}{\partial W_{\mathrm{col}}}=X^TG.
\]
The two expressions differ only by the convention for storing the weight matrix.

\section{Outer-product expression}

Let
\[
  \frac{\partial y}{\partial z}
\]
be treated as a column vector for elementwise activation, with entries $\sigma'(z_i)$.  Then the derivative of $y$ with respect to $W$ can be represented row by row as
\[
  \frac{\partial y}{\partial W}\sim
  \begin{bmatrix}
  \frac{\partial y_1}{\partial z_1}x^T\\
  \frac{\partial y_2}{\partial z_2}x^T\\
  \vdots\\
  \frac{\partial y_m}{\partial z_m}x^T
  \end{bmatrix}
  =\left(\frac{\partial y}{\partial z}\right)x^T.
\]
For the scalar loss, replace $\partial y/\partial z$ by the backpropagated vector $\delta=\partial L/\partial z$.

\section{Numerical example}

The note gives a simple example with
\[
  m=2,\qquad n=3,\qquad x=[1,2,3]^T,\qquad z=[0.1,-1.0]^T.
\]
For sigmoid,
\[
  \sigma'(0.1)\approx0.2494,\qquad
  \sigma'(-1.0)\approx0.1966.
\]
Therefore
\[
  \frac{\partial y}{\partial W}\sim
  \begin{bmatrix}0.2494\\0.1966\end{bmatrix}
  \begin{bmatrix}1&2&3\end{bmatrix}
  =
  \begin{bmatrix}
  0.2494&0.4988&0.7482\\
  0.1966&0.3932&0.5898
  \end{bmatrix}.
\]
Each row is the input vector scaled by the derivative of the corresponding output.

\section{Einstein notation proof}

An index calculation makes the chain rule explicit.  Let $y=f(z)$ and $z=Wx$.  For a matrix entry $A_{pq}=W_{pq}$,
\[
  \frac{\partial z_i}{\partial A_{pq}}
  =\frac{\partial}{\partial A_{pq}}
  \sum_j A_{ij}x_j
  =\delta_{ip}x_q.
\]
Then
\[
  \frac{\partial L}{\partial A_{pq} }
  =\sum_i \frac{\partial L}{\partial y_i}
  \frac{\partial y_i}{\partial z_p}x_q.
\]
For elementwise activation, this becomes
\[
  \frac{\partial L}{\partial A_{pq}}
  =\frac{\partial L}{\partial z_p}x_q.
\]
This is the component form of $\partial L/\partial W=\delta x^T$.

\section{Kronecker product viewpoint}

The note also explores vectorization.  If $y=f(Wx)$ and we vectorize $W$, then the Jacobian with respect to $\operatorname{vec}(W)$ can be written using Kronecker products.  The important identity is that changing one matrix entry affects only one row of $z$, which is encoded by a Kronecker delta.

In batch form, if $X$ is a matrix of inputs and $\Delta$ is the matrix of upstream gradients, then the gradient becomes
\[
  \frac{\partial L}{\partial W}=\Delta X^T.
\]
This is the same outer-product rule, summed over the batch dimension.

\section{Softmax--cross-entropy simplification}

For the final layer with softmax and cross entropy, the derivative simplifies to
\[
  \frac{\partial L}{\partial z}=\hat y-t.
\]
Therefore the final-layer weight gradient is
\[
  \frac{\partial L}{\partial W}=(\hat y-t)x^T.
\]
This is why implementation code often looks much simpler than the theoretical tensor derivation.

\section{A wider interpretation}

The key lesson is that matrix calculus is compressed tensor calculus.  The full derivative of a vector with respect to a matrix is a higher-order object, but backpropagation contracts it immediately with an upstream gradient.  The result is an outer product: error signal times input.  This is the local rule behind every dense neural-network layer.

\section{Backpropagation as cotangent pullback}
For a layer \(z=Wx+b\), the forward map sends activations forward, while backpropagation sends covectors backward:
\[
\bar x=W^T\bar z.
\]
This is the pullback on cotangent vectors.
\section{Why the three-dimensional tensor collapses}
The derivative of \(z_i=\sum_j W_{ij}x_j+b_i\) with respect to \(W_{ab}\) is
\[
\frac{\partial z_i}{\partial W_{ab}}=\delta_{ia}x_b.
\]
Contracting with upstream gradient \(\bar z_i\) gives
\[
\frac{\partial L}{\partial W_{ab}}=\bar z_a x_b.
\]
Thus the huge derivative tensor disappears after contraction, leaving an outer product.
\section{Batch form}
For a batch, gradients add over examples:
\[
\nabla_W L=\bar Z X^T.
\]
This is matrix multiplication as accumulated outer products.
\section{Backpropagation as sparse tensor contraction}
Backpropagation is not magic tensor cancellation.  It is cotangent pullback plus contraction of sparse Jacobians, and the familiar weight gradient is the outer product of upstream error with input activation.
